\documentclass[12pt,letterpaper]{article}

% ── Packages ─────────────────────────────────────────────────────────────────
\usepackage[margin=1in]{geometry}
\usepackage{setspace}
\usepackage{titlesec}
\usepackage{booktabs}
\usepackage{tabularx}
\usepackage{array}
\usepackage{longtable}
\usepackage{graphicx}
\usepackage{xcolor}
\usepackage{hyperref}
\usepackage{enumitem}
\usepackage{parskip}
\usepackage{microtype}
\usepackage{fancyhdr}
\usepackage{amsmath}
\usepackage{caption}

% ── Color scheme ─────────────────────────────────────────────────────────────
\definecolor{accentgreen}{RGB}{42,106,42}
\definecolor{lightgray}{RGB}{245,245,245}
\definecolor{darkgray}{RGB}{80,80,80}

% ── Hyperref setup ───────────────────────────────────────────────────────────
\hypersetup{
  colorlinks=true,
  linkcolor=accentgreen,
  urlcolor=accentgreen,
  citecolor=accentgreen
}

% ── Section formatting ───────────────────────────────────────────────────────
\titleformat{\section}{\large\bfseries\color{accentgreen}}{{\thesection.}}{0.5em}{}[\titlerule]
\titleformat{\subsection}{\normalsize\bfseries\color{darkgray}}{{\thesubsection}}{0.5em}{}
\titlespacing*{\section}{0pt}{1.2em}{0.5em}
\titlespacing*{\subsection}{0pt}{0.8em}{0.3em}

% ── Header / footer ──────────────────────────────────────────────────────────
\pagestyle{fancy}
\fancyhf{}
\fancyhead[L]{\small\color{darkgray}NourishAI — Project Report}
\fancyhead[R]{\small\color{darkgray}Entrepreneurial Mindset}
\fancyfoot[C]{\small\thepage}
\renewcommand{\headrulewidth}{0.4pt}

% ── Column types ─────────────────────────────────────────────────────────────
\newcolumntype{L}[1]{>{\raggedright\arraybackslash}p{#1}}
\newcolumntype{C}[1]{>{\centering\arraybackslash}p{#1}}
\newcolumntype{R}[1]{>{\raggedleft\arraybackslash}p{#1}}

\setlength{\parskip}{6pt}
\onehalfspacing

% ─────────────────────────────────────────────────────────────────────────────
\begin{document}

% ── Title Page ───────────────────────────────────────────────────────────────
\begin{titlepage}
  \centering
  \vspace*{2cm}
  {\Huge\bfseries\color{accentgreen} NourishAI\par}
  \vspace{0.4cm}
  {\Large\bfseries An AI-Powered RAG Meal Planning Platform\par}
  \vspace{0.6cm}
  {\large\color{darkgray} Project Report\par}
  \vspace{2cm}

  \rule{0.5\textwidth}{0.4pt}\par
  \vspace{0.4cm}
  {\normalsize
    Jose Rafael Yaptangco \quad Vaibhavi Dalvi \\[4pt]
    Sahil Kassam \quad Sunny Bhargava
  \par}
  \vspace{0.4cm}
  \rule{0.5\textwidth}{0.4pt}\par

  \vspace{1.5cm}
  {\normalsize Fordham University\\[4pt]
  Entrepreneurial Mindset\par}

  \vspace{1cm}
  {\normalsize May 2026\par}

  \vfill
  \begin{abstract}
    \noindent
    NourishAI is a full-stack, AI-driven meal planning web application built on a
    Retrieval-Augmented Generation (RAG) architecture. The platform allows users to
    discover recipes through natural-language conversation, assemble personalized
    weekly schedules, generate consolidated grocery lists, and export meal prep guides.
    This report covers the application's purpose and technical design, the market
    context and business strategy drawn from our operations and marketing plans, and a
    rigorous empirical evaluation of the RAG retrieval pipeline across five endpoints
    using 60 LLM-judged test queries. Mean Precision@5 ranges from 0.633 to 0.783
    across endpoints, and mean Recall@20 reaches up to 0.343, demonstrating strong
    top-of-list relevance against a challenging, small-index backdrop of 697 recipes.
  \end{abstract}
\end{titlepage}

\newpage
\tableofcontents
\newpage

% ─────────────────────────────────────────────────────────────────────────────
\section{Introduction}

The decision of what to eat each week is, for most people, a repetitive administrative burden
rather than a source of joy. Industry data bear this out: early adopters of AI meal-planning
tools save an average of three hours per week on meal planning and reduce household food
waste by up to 40\%. Yet the tools available today force users to choose between static,
form-driven apps that offer structure without intelligence, and general-purpose AI assistants
that offer intelligence without structure or memory.

NourishAI was built to close that gap. The project began as an exploration of how
Retrieval-Augmented Generation (RAG) could be applied in a domain where the user's intent
is inherently personal, iterative, and hard to anticipate: everyday cooking. Rather than
training a bespoke model, the team assembled a curated vector index of recipes and connected
it to a large language model (Gemini 2.5 Flash via LiteLLM) through a FastAPI backend.
The result is a platform that understands what a user means when they type ``something
quick and comforting for a busy Tuesday'' and returns semantically relevant recipes
instantly.

This report is organized as follows. Section~\ref{sec:app} describes the application's
features and architecture. Section~\ref{sec:market} situates it in the competitive landscape.
Section~\ref{sec:ops} summarizes the operational plan. Section~\ref{sec:eval-method}
defines the evaluation methodology. Section~\ref{sec:results} presents quantitative results,
and Section~\ref{sec:discussion} interprets them. Section~\ref{sec:conclusion} concludes.

% ─────────────────────────────────────────────────────────────────────────────
\section{Application Overview}
\label{sec:app}

\subsection{Core Features}

NourishAI exposes four main surfaces to the user:

\begin{enumerate}[leftmargin=*, itemsep=2pt]
  \item \textbf{Conversational Chat.} A natural-language interface powered by Gemini 2.5
    Flash retrieves up to 20 semantically similar recipes from the Pinecone vector index
    and presents them with images, ingredients, and step-by-step instructions. The LLM
    can suggest full weekly plans directly from the chat view, which are applied to the
    planner with a single click.

  \item \textbf{Weekly Planner.} A 7-day $\times$ 3-meal interactive grid. Users fill cells
    manually, by AI prompt (``arrange a high-protein week with Italian dinners''), or via
    one-click Magic Auto-Fill that derives a preference query from the user's saved
    favorites. Nutritional totals (calories, protein, carbs, fat) update in real time as
    recipes are added or removed. Completed weeks can be saved as recurring templates.

  \item \textbf{Grocery List \& Meal Prep Guide.} The grocery endpoint consolidates
    ingredients from all meals in the current week's plan, deduplicates quantities, and
    organizes them by category. The prep guide endpoint produces a step-by-step
    batch-prep schedule ordered for kitchen efficiency. Both can be copied to the clipboard
    or exported as CSV.

  \item \textbf{Google Calendar Integration.} Users may push their weekly plan to Google
    Calendar via OAuth~2.0, creating one event per meal with the full ingredient list in
    the event description.
\end{enumerate}

\subsection{Technical Architecture}

Table~\ref{tab:stack} summarizes the technology stack. The backend is a single FastAPI
application deployed as a Vercel serverless function. The retrieval layer is Pinecone
(serverless index, \texttt{text-embedding-3-small}, 1536 dimensions). User data and
favorites are persisted in SQLite for local development and Supabase Postgres in production.

\begin{table}[h!]
  \centering
  \caption{Core technology stack}
  \label{tab:stack}
  \begin{tabularx}{\textwidth}{L{3.2cm} X}
    \toprule
    \textbf{Layer} & \textbf{Technology} \\
    \midrule
    Frontend      & React 18 + Vite, Vanilla CSS (dark-mode theme) \\
    Backend       & Python 3.13, FastAPI, Uvicorn \\
    LLM           & Gemini 2.5 Flash via LiteLLM (\texttt{gemini/gemini-2.5-flash}) \\
    Embeddings    & OpenAI \texttt{text-embedding-3-small} (1536 dims) \\
    Vector Store  & Pinecone Serverless (697 recipe vectors at evaluation time) \\
    Relational DB & SQLite (local) / Supabase Postgres (production) \\
    Auth          & SHA-256 password hashing; Google OAuth 2.0 for Calendar \\
    Deployment    & Vercel (frontend + API serverless functions) \\
    \bottomrule
  \end{tabularx}
\end{table}

\subsection{RAG Pipeline Design}

Every user query follows the same retrieval path:

\begin{enumerate}[itemsep=2pt]
  \item The raw query string is embedded with \texttt{text-embedding-3-small}.
  \item Pinecone returns the top-$k$ nearest neighbours (configurable per endpoint:
    $k \in \{20, 40, 50\}$).
  \item The matched recipe metadata (title, category, cuisine, ingredients, nutrition,
    instructions) is injected into a structured LLM prompt.
  \item Gemini 2.5 Flash generates a response constrained to the retrieved context,
    ensuring factual grounding and preventing hallucination of non-existent recipes.
  \item Structured outputs (meal plans, grocery lists, prep guides) are validated by
    Pydantic schemas before being returned to the frontend.
\end{enumerate}

The recipe index was populated from three sources: TheMealDB (public API), a curated
synthetic generation pass using Gemini, and Spoonacular/RecipeNLG data. Each recipe
is stored as a Pinecone vector with metadata fields for title, category, cuisine, ingredients
(JSON), instructions, and nutrition (calories, protein, carbs, fat).

% ─────────────────────────────────────────────────────────────────────────────
\section{Market Context \& Business Strategy}
\label{sec:market}

\subsection{Market Opportunity}

The AI-driven meal planning app market was valued at \$972 million in 2024 and is
projected to reach \$11.6 billion by 2034, growing at a 28.1\% CAGR. North America leads
with 32.1\% market share (\$312 million in 2024), and the U.S.\ alone contributed
\$271.5 million. Adoption rates jumped 47\% in 2025, with early adopters saving an average
of three hours per week on meal planning and reducing food waste by up to 40\%.

Individual consumers represent 72\% of the user base, with the 30--49 demographic
accounting for 51\% of active users. Subscription revenue delivers 44.89\% of total market
revenue, while B2B enterprise licensing is expanding at 15.37\% CAGR.

\subsection{Competitive Landscape}

The market is fragmented but consolidating around five major players (Mealime, Paprika,
Forks Over Knives, MealPrepPro, and Mealboard) that collectively hold 48\% of market
share. NourishAI enters this space occupying a strategically distinct position: the only
platform that combines conversational AI input, persistent user profiles, end-to-end workflow
completion (intent $\to$ plan $\to$ grocery list $\to$ prep guide), and accessible pricing.

Table~\ref{tab:competitive} maps key competitors across the dimensions most relevant to
target users.

\begin{table}[h!]
  \centering
  \caption{Competitive feature comparison}
  \label{tab:competitive}
  \small
  \begin{tabularx}{\textwidth}{L{2.8cm} C{1.5cm} C{1.5cm} C{1.5cm} C{1.5cm} C{2cm}}
    \toprule
    \textbf{Feature} & \textbf{NourishAI} & \textbf{Mealime} & \textbf{Noom} & \textbf{PlateJoy} & \textbf{HelloFresh} \\
    \midrule
    Conversational AI      & \checkmark & -- & -- & -- & -- \\
    Persistent Profile     & \checkmark & -- & Partial & Static & -- \\
    Auto Grocery List      & \checkmark & \checkmark & -- & \checkmark & \checkmark \\
    Grocery Delivery       & Roadmap    & Partial & -- & Instacart & Core \\
    Full Workflow          & \checkmark & Partial & -- & Partial & \checkmark \\
    Diet Flexibility       & Any        & Limited & Cal-density & \checkmark & Limited \\
    Monthly Price          & \$7--\$15  & Free/\$2.99 & \$17--\$70 & \$8--\$12 & \$60--\$100+ \\
    \bottomrule
  \end{tabularx}
\end{table}

Three strategic gaps exist in the current market that NourishAI is positioned to capture:

\begin{itemize}[itemsep=2pt]
  \item \textbf{The Persistent AI Gap.} No competitor combines natural-language input
    with persistent memory and an integrated downstream workflow. ChatGPT has the
    language capability but not the integration or memory; Mealime has the workflow
    but not conversational AI. NourishAI bridges both.

  \item \textbf{The Family Planning Gap.} The 30--49 demographic represents 51\% of
    users, yet most competitors optimize for individuals. Family-tier planning---managing
    preferences across household members, scaling recipes, consolidating multi-person
    grocery lists---is underbuilt across all major competitors.

  \item \textbf{The Meal Kit Churn Pool.} Meal kit subscribers (HelloFresh, Blue Apron)
    have already demonstrated willingness to pay for meal-planning convenience, then
    churned due to rigidity and cost. NourishAI offers equivalent structure at grocery
    store prices with full dietary freedom.
\end{itemize}

\subsection{Value Proposition and Positioning}

NourishAI's core positioning statement is: \textit{the personal digital sous-chef that turns
what you say into a week of real meals, tailored exactly to your life, not a template of it.}

The platform is deliberately framed around \textit{wellness}, not ``health.'' Health as a
marketing term carries clinical baggage that excludes the casual user who simply wants to
stop ordering takeout four nights a week. Wellness is broader and more inclusive: it
serves the fitness enthusiast tracking macros just as effectively as the parent who types
``something quick for a weeknight with whatever's in the fridge.''

The competitive moat is not the AI model itself (which any competitor can access) but the
compounding behavioral data layer. As the platform accumulates each user's preferences,
favorites, and feedback patterns, the switching cost rises through \textit{earned utility}
rather than lock-in mechanisms.

% ─────────────────────────────────────────────────────────────────────────────
\section{Operations \& Development Plan}
\label{sec:ops}

\subsection{Development Stages}

The platform was built in six sequential stages mirroring the operations plan:

\begin{enumerate}[itemsep=2pt, leftmargin=*]
  \item \textbf{Recipe Database Curation.} A corpus of recipes was sourced from
    TheMealDB, Spoonacular, and synthetic generation. Each entry was tagged with
    cuisine, category, ingredients, instructions, and nutritional metadata.

  \item \textbf{Data Processing \& Vector Embedding.} Recipes were embedded using
    OpenAI's \texttt{text-embedding-3-small} model (1536 dimensions), capturing
    semantic meaning for natural-language retrieval.

  \item \textbf{Vector Database Setup.} Embeddings were loaded into Pinecone
    Serverless, supporting fast ANN search across 697 indexed vectors at evaluation time.

  \item \textbf{RAG Pipeline \& LLM Integration.} The retrieval layer was connected to
    Gemini 2.5 Flash via LiteLLM. All structured outputs are validated by Pydantic schemas
    before delivery to the frontend.

  \item \textbf{User Profile \& Persistence.} A favorites system and persistent schedule
    storage (SQLite / Supabase) enable the platform to serve returning users with
    preference-aware recommendations without re-entering baseline information.

  \item \textbf{Frontend Development.} A React 18 + Vite SPA provides the conversational
    chat interface, planner grid, grocery and prep workflows, and Google Calendar export.
\end{enumerate}

\subsection{Infrastructure \& Cost Profile}

All vendor relationships are API-based SaaS subscriptions with no physical supply chain.
The estimated operating cost at MVP scale is approximately \$200--\$500 per month,
comprising LLM API calls (\$0.01--\$0.05 per user session), Pinecone (\$70/month at
production tier), cloud hosting on Vercel (free tier for MVP), and recipe data APIs.
Stripe processes subscription payments at 2.9\% + \$0.30 per transaction.

\subsection{Go-to-Market Timeline}

Key milestones from the operations plan:

\begin{center}
\small
\begin{tabular}{L{8cm} C{3cm}}
  \toprule
  \textbf{Milestone} & \textbf{Target} \\
  \midrule
  MVP live (web); closed beta (50--100 users)   & Month 5  \\
  Public launch (web + mobile)                   & Month 8  \\
  Grocery API integration (Kroger / Instacart)   & Month 10 \\
  1,000 monthly active users                     & Month 12 \\
  Premium tier launch (family plans)             & Month 14 \\
  5,000 MAU; first B2B wellness pilot            & Month 18 \\
  Seed round target (\$500K--\$1M)               & Month 18--24 \\
  15,000 monthly active users                    & Month 24 \\
  \bottomrule
\end{tabular}
\end{center}

% ─────────────────────────────────────────────────────────────────────────────
\section{RAG Evaluation Methodology}
\label{sec:eval-method}

\subsection{Motivation}

Retrieval quality is the critical dependency in any RAG system: if the wrong recipes are
surfaced, no amount of LLM quality can rescue a coherent, relevant response. We designed
a rigorous offline evaluation covering all five RAG-powered endpoints to measure how
well the Pinecone index satisfies diverse, realistic user queries.

\subsection{Endpoints Evaluated}

Five endpoints were evaluated, each with a distinct top-$k$ budget reflecting production
configuration:

\begin{center}
\small
\begin{tabular}{L{4.5cm} C{2cm} L{6cm}}
  \toprule
  \textbf{Endpoint} & \textbf{top-$k$} & \textbf{Query style} \\
  \midrule
  \texttt{/api/recipes/search}    & 20 & Short keyword / ingredient fragments \\
  \texttt{/api/chat}              & 20 & Conversational, situational requests \\
  \texttt{/api/planner/prompt}    & 50 & Multi-day planning prompts \\
  \texttt{/api/planner/autofill}  & 40 & Compact preference profiles \\
  \texttt{/api/recommendations}   & 40 & Personal taste expressions \\
  \bottomrule
\end{tabular}
\end{center}

\subsection{Test Case Generation}

Twelve synthetic test queries per endpoint (60 total) were generated by Gemini 2.5 Flash
using a controlled prompt enforcing diversity across query length, dietary constraints,
occasions, cuisines, and persona types. Generation rules prohibited food-blog language
(``delicious'', ``perfect''), required natural informal phrasing, and mandated that no two
queries share an opening word.

\subsection{Relevance Judging}

Relevance was determined by an LLM judge (Gemini 2.5 Flash) rather than keyword
heuristics, enabling nuanced assessment of semantic intent. For each query, the judge
evaluated whether each retrieved recipe was relevant, applying an inclusive standard:
same cuisine, matching category, or a key ingredient all count as relevant.

The \textbf{reference pool} was constructed by retrieving the top-100 results per query from
Pinecone, approximating the universe of relevant items in the index. This pool serves as
the denominator for recall.

\subsection{Metrics}

For each query $q$, retrieved set $R_k$ (top-$k$ results), and reference pool $P$:

\[
\text{Precision@}k = \frac{|\{r \in R_k : \text{relevant}(r)\}|}{k}
\qquad
\text{Recall@}k = \frac{|\{r \in R_k : \text{relevant}(r)\}|}{|\{r \in P : \text{relevant}(r)\}|}
\]

$K$ values evaluated: 1, 3, 5, 10, and 20. Metrics are averaged across the 12 queries
per endpoint to produce endpoint-level summaries.

% ─────────────────────────────────────────────────────────────────────────────
\section{Evaluation Results}
\label{sec:results}

\subsection{Summary Table}

Table~\ref{tab:summary} presents mean Precision@$K$ and Recall@$K$ for each endpoint
across all $K$ values, from the primary evaluation run (Run ID: 20260425\_124536,
697 indexed vectors).

\begin{table}[h!]
  \centering
  \caption{Mean Precision@K and Recall@K by endpoint (12 queries each, $n=60$ total)}
  \label{tab:summary}
  \small
  \begin{tabularx}{\textwidth}{L{3.6cm} C{1.1cm} C{1.1cm} C{1.1cm} C{1.1cm} C{1.1cm} C{1.1cm} C{1.1cm} C{1.1cm} C{1.1cm} C{1.1cm}}
    \toprule
    & \multicolumn{5}{c}{\textbf{Precision@K}} & \multicolumn{5}{c}{\textbf{Recall@K}} \\
    \cmidrule(lr){2-6}\cmidrule(lr){7-11}
    \textbf{Endpoint} & @1 & @3 & @5 & @10 & @20 & @1 & @3 & @5 & @10 & @20 \\
    \midrule
    /api/recipes/search    & .917 & .806 & .717 & .633 & .562 & .026 & .063 & .092 & .157 & .275 \\
    /api/chat              & .833 & .750 & .717 & .658 & .617 & .030 & .076 & .107 & .152 & .257 \\
    /api/planner/prompt    & .917 & .806 & .783 & .617 & .596 & .022 & .052 & .082 & .118 & .213 \\
    /api/planner/autofill  & .833 & .778 & .767 & .750 & .704 & .025 & .062 & .105 & .195 & .343 \\
    /api/recommendations   & .583 & .583 & .633 & .592 & .588 & .020 & .059 & .106 & .155 & .257 \\
    \midrule
    \textbf{Overall mean}  & .817 & .745 & .723 & .650 & .613 & .025 & .062 & .098 & .155 & .269 \\
    \bottomrule
  \end{tabularx}
\end{table}

\subsection{Precision Analysis}

\textbf{Top-of-list precision is consistently high.} Precision@1 exceeds 0.80 for four of
five endpoints, with search and planner/prompt achieving 0.917---meaning the very first
retrieved result is relevant nearly every time. This is the most practically important metric:
users form their first impression from the top card displayed.

\textbf{Precision degrades gracefully with $K$.} The decline from P@1 to P@20 ranges
from 0.245 points (planner/prompt: 0.917 $\to$ 0.596) to 0.129 points (recommendations:
0.583 $\to$ 0.588). The flat profile for the recommendations endpoint reflects the broader
top-40 retrieval budget absorbing more marginal results.

\textbf{Autofill demonstrates the most stable precision curve.} P@5 = 0.767 and
P@10 = 0.750 indicate that the preference-profile query string (derived from user favorites)
yields a densely relevant result set across the full top-10.

\textbf{Recommendations shows lower absolute precision} (P@5 = 0.633) but remains above
chance. Several failure cases involved queries that were intentionally narrow (``zucchini
and feta,'' ``no dairy dessert'') in a 697-vector index with limited coverage of those
exact subcategories, indicating an index size limitation rather than an embedding quality
failure.

\subsection{Recall Analysis}

\textbf{Recall values are bounded by index size.} With only 697 total vectors, many
queries have a large relevant fraction of the index (e.g., ``quick chicken thing'' has 98
relevant recipes in the reference pool of 100---virtually the entire index is relevant). In
such cases, even retrieving 20 results can only recover $\approx$20\% of the relevant set.
This structural ceiling makes recall an informative but bounded metric at current scale.

\textbf{Autofill achieves the best recall}, reaching R@20 = 0.343---meaning that the
40-result retrieval budget recovers over a third of all relevant index items on average.

\textbf{Chat recall is competitive with search} despite more conversational queries
(R@20 = 0.257 vs.\ 0.275), validating that the embedding model handles situational
natural-language input as effectively as short keyword queries.

\textbf{Planner/prompt recall is lower} (R@20 = 0.213), expected given the complexity of
multi-constraint planning prompts that may embed several cuisines, dietary restrictions, and
time constraints in a single query---a harder retrieval problem than a focused dish search.

\subsection{Notable Query-Level Observations}

Several queries revealed specific coverage gaps:

\begin{itemize}[itemsep=2pt]
  \item ``No dairy creamy pasta sauce'' (search): P@5 = 0.00. The index contained only
    5 relevant results in the reference pool, suggesting sparse coverage of dairy-free
    pasta variations.

  \item ``What's a good side dish for grilled salmon?'' (chat): P@5 = 0.00. Side dish
    framing (rather than a main dish query) does not match the recipe metadata schema,
    which categorizes by entree type rather than meal role.

  \item ``Something Thai'' (search): P@5 = 1.00, R@20 = 0.556---a dense cluster of
    Thai recipes was retrieved perfectly, demonstrating the embedding model's strength
    at cuisine-level retrieval.

  \item ``Baking with apples'' (recommendations): R@20 = 0.875---the highest recall of
    any single query. A very small reference pool (8 relevant items) was almost entirely
    recovered in the top-40 results.
\end{itemize}

% ─────────────────────────────────────────────────────────────────────────────
\section{Discussion}
\label{sec:discussion}

\subsection{Interpretation of Results}

The evaluation demonstrates that the RAG retrieval pipeline performs well on its primary
task: surfacing relevant recipes at the top of the ranked list. Mean Precision@5 of 0.723
across all endpoints means that, on average, more than 3.5 of the first 5 results shown to
the user are relevant to their query---a threshold sufficient to generate a satisfying user
experience.

The relatively low recall values are a direct artifact of the small index (697 vectors) rather
than a retrieval quality failure. At production scale with a larger recipe corpus, the recall
denominator would be the set of truly relevant recipes in a much larger index, and the
absolute recall numbers would become more meaningful. The current evaluation is best
interpreted as \textit{precision-centric}: the system is good at finding \textit{some}
good results, but the index is not yet large enough to demonstrate comprehensive coverage.

\subsection{Limitations}

\begin{itemize}[itemsep=2pt]
  \item \textbf{Index size.} 697 vectors is small for a production recipe application.
    The operations plan targets 200--300 recipes for MVP with ongoing weekly additions;
    the current index slightly exceeds that but remains below the coverage needed for
    niche query categories.

  \item \textbf{LLM judge calibration.} Relevance judgments were made by the same
    model family (Gemini) used for query generation. While the judge prompt was
    carefully crafted to be inclusive, inter-rater agreement with human judges was not
    measured.

  \item \textbf{Single embedding model.} All queries and documents use
    \texttt{text-embedding-3-small}. Hybrid retrieval (sparse + dense) or fine-tuned
    food-domain embeddings could improve performance on niche or ingredient-specific
    queries.
\end{itemize}

\subsection{Business Implications}

From a product perspective, the evaluation results support the core value proposition.
High top-of-list precision means that the first recipe a user sees is almost always relevant,
reducing the friction of sifting through poor results. The planner/autofill endpoint's stable
precision curve (P@5 through P@10 $\approx$ 0.75--0.77) is particularly important: the
one-click auto-fill feature draws from a top-40 pool, and this evaluation shows that pool
is high quality. Expanding the recipe index is the clearest path to improving recall and
supporting a wider range of dietary preferences and cuisines---a direct input to the weekly
database update cadence described in the operations plan.

% ─────────────────────────────────────────────────────────────────────────────
\section{Conclusion}
\label{sec:conclusion}

NourishAI demonstrates that a RAG architecture is an effective and practical foundation
for a conversational meal planning application. The platform successfully combines
semantic vector retrieval with a capable LLM to deliver an experience that competitors
built on static recipe databases or form-driven filters cannot replicate at the same price
point and with the same breadth of natural-language input support.

The quantitative evaluation across 60 LLM-judged test queries confirms strong retrieval
precision---averaging 0.723 at $K=5$ across five endpoints---with recall appropriately
bounded by the current index size. These results validate the architectural choices made
during development and identify a clear growth path: expanding the recipe corpus, adding
hybrid retrieval, and incorporating user behavioral feedback into re-ranking.

From a business standpoint, NourishAI enters a \$972 million market growing at 28.1\%
CAGR with a differentiated positioning: the only platform that combines conversational AI,
persistent user memory, and end-to-end workflow completeness at a freemium price point
accessible to casual wellness users and fitness enthusiasts alike. The technical
foundation is proven; the market opportunity is established; the next step is scale.

\vspace{1em}
\noindent\rule{\textwidth}{0.4pt}
\vspace{0.5em}

\noindent\textit{NourishAI was developed as part of the Entrepreneurial Mindset course at
Fordham University by Jose Rafael Yaptangco, Vaibhavi Dalvi, Sahil Kassam, and Sunny
Bhargava (Spring 2026).}

\end{document}
